% Plantilla de analysis.pdf (Quiz 03). 
% Los comentarios % bajo cada apartado son pistas: no salen en el PDF.
% Limite de la entrega: 4 paginas.
% Compilar: pdflatex analysis_plantilla.tex
\documentclass[11pt,a4paper]{article}
\usepackage[utf8]{inputenc}
\usepackage[T1]{fontenc}
\usepackage{lmodern}
\usepackage[spanish,es-noshorthands]{babel}
\usepackage{amsmath,amssymb}
\usepackage{booktabs,array}
\usepackage{enumitem}
\usepackage{titlesec}
\usepackage[a4paper,top=1.4cm,bottom=1.5cm,left=1.7cm,right=1.7cm]{geometry}

\setlength{\parindent}{0pt}
\setlength{\parskip}{2pt}
\setlength{\emergencystretch}{2em}
\setlist{nosep,leftmargin=1.3em}
\titleformat{\section}{\large\bfseries}{}{0pt}{}
\titlespacing*{\section}{0pt}{9pt}{2pt}
\newcolumntype{L}[1]{>{\raggedright\arraybackslash}p{#1}}
\newcommand{\phys}{\operatorname{phys}}
% Apartado: \sub{a}{Titulo corto}
\newcommand{\sub}[2]{\medskip\par\noindent\textbf{#1) #2.}\ }

\begin{document}

\begin{center}
{\Large\bfseries Estructuras de Datos --- Quiz 03}\\[1pt]
{\large Dynamic Arrays and Amortized Analysis} \quad$\cdot$\quad analysis.pdf \quad$\cdot$\quad Semestre 2026-2
\end{center}
\vspace{-6pt}
\textbf{Nombre:} Kevin Javier Gonzalez \hfill \textbf{Documento:} 1093432420 \hfill \textbf{Grupo:} 3

\section*{Notación y modelo de costo}
% k, C, n (incluye transitorio k+1), m, g<=m.
% Modelo: escribir=1, copiar en resize=1, aritmetica O(1).


%=====================================================================
\section*{P1 --- Arreglo dinámico desde cero}

\sub{a}{Implementación}
% Referencia a src/DynamicArray.java. Representacion: data, size, capacity.
En el archivo de DynamicArray.Java esta el Codigo.

\sub{b}{Invariantes}
% Minimo 3, validos antes y despues de cada operacion publica.
% Sugeridos: 0<=size<=capacity; capacity=data.length>=1; data[0..size-1] validos;
% data[size..] = null; capacity solo se duplica cuando size=capacity.
\begin{enumerate}[label=\textbf{I\arabic*.}]
\item 0 <= size <= capacity
\item capacity >= 1
\item valores no nulos en data estan desde [0, size - 1], data[size] == null
\end{enumerate}
% Justificar la preservacion en append / removeLast / resize.


\sub{c}{Complejidad worst-case}\par\smallskip
\renewcommand{\arraystretch}{2}
\begin{tabular}{@{}L{0.16\textwidth}L{0.2\textwidth}L{0.55\textwidth}@{}}
\toprule
Operación & Worst-case & Justificación\\
\midrule
\texttt{get}       & $\Theta(1)$  & es simplemente una busqueda, el array por defecto ya lo hace en O(1) \\ 
\texttt{set}        & $\Theta(1)$ & igualmente al set, es una busqueda de costo 1 + una asignacion de costo 1 en todo los casos\\
\texttt{append}     & O(n) & en el peor caso tengo que construir un nuevo array y dato por dato del anterior asignarlo en la misma pos del nuevo array, claro n es size dado que mas alla de size no toca asignar nada\\
\texttt{removeLast} & $\Theta(1)$ & aca igualmente es siempre el mismo valor, hago una busqueda de costo 1, y hago una asignacion de costo 1, y luego simplemente hago size-- que tambien es de costo 1 \\
\bottomrule
\end{tabular}
\renewcommand{\arraystretch}{1}


\sub{d}{Worst-case vs.\ amortizado}
% Por que el append caro no contradice el amortizado constante.
% Resize solo con size=C, cuesta C, deja 2C. Sumar los resizes de N appends.

cuando hago append si $size$ < a $capacity$ entonces solo tengo que hacer una operacion de asignacion de coste 1, entonces cada que se duplica la capacity tengo $\frac{capacity}{2}$ veces para hacer operaciones de costo 1

%=====================================================================
\section*{P2 --- Snake sobre \texttt{DynamicArray} con desplazamiento}

\sub{a}{Costo de \texttt{removeTail}}
% Iteraciones del bucle cuando size = k+1.

cuando se llama a remove tail en esa implementacion, toca eliminar la cola y luego, reasignar los valores que quedan una posicion hacia atras, el costo seria O(n -1) -> O(n), con n siendo el size.

\sub{b}{Costo total de $m$ movimientos}
% Longitud constante k. Sumar append + desplazamiento (+ resize inicial).

siendo de longitud $k$, hacer $m$ movimientos seria que en cada movimiento tengo el costo $O(k)$ de mover la serpiente, m veces esto, osea el costo final seria $O(k \cdot m)$

\sub{c}{Validez de la afirmación} % por hacer
% Veredicto + prueba formal. Costo(MOVE) = append + RemoveTail.
% Comparar sum a_i con sum r_i.

esto es falso, $move()$ no es amortizado porque en cada operacion de $move()$ tengo que hacer n (size) actualizaciones de posicion de la serpiente lo que hace que en promedio sea $\frac{O(n \cdot n)}{n} \rightarrow O(n)$ veces el costo amortizado, a diferencia del append, que en n operaciones el coste es $\frac{O(2n) + O(n)}{n} \rightarrow O(3) \rightarrow O(1)$


\sub{d}{Bottleneck} % por hacer
% Que fija la representacion y por que obliga a mover k elementos.

el bottleneck seria que como en cado movimiento que no coma tengo que hacer $eliminarCola()$ entonces tendria n para ajustar el array dado que la nueva cola debe estar en la posicion 0, pero a su vez como esta en movimiento tengo que hacer n operaciones de calcular su nueva posicion, esto haria que para m operaciones sin comer tendria (n+n)*m operaciones

%=====================================================================
\section*{P3 --- Rediseño: representación interna}

\sub{a}{Variables de estado}
\begin{itemize}
\item \texttt{data}:
\item \texttt{start}:
\item \texttt{size}:
\item \texttt{capacity}: 
\end{itemize}

\sub{b}{Mapeo lógico $\to$ físico}
\[
\phys(i)=
\]
% Demostrar: (1) 0<=phys(i)<C. (2) inyectiva. (3) la celda phys(size) esta libre si size<C.


\sub{c}{Operaciones}
% Codigo en src/Snake.java. Costo de cada una.
\begin{itemize}
\item \texttt{addHead}:
\item \texttt{removeTail}:
\item \texttt{get}:
\end{itemize}

\sub{d}{\texttt{resize}}
% Como conserva el orden logico aunque la cola este a la derecha de la cabeza.


\sub{e}{Invariantes y corrección}
\begin{enumerate}[label=\textbf{S\arabic*.}]
\item
\item
\item
\item
\end{enumerate}
% Preservacion en addHead / removeTail / resize.


\sub{f}{Traza}
% Estado tras cada movimiento (size despues de removeTail). Celdas vacias con ``$-$''.
\par\smallskip
\renewcommand{\arraystretch}{1.9}
\begin{tabular}{@{}L{2.3cm}L{4.6cm}c c c L{1.9cm}L{3cm}@{}}
\toprule
Mov. & Arreglo físico & \texttt{start} & \texttt{size} & \texttt{capacity} & ¿resize? & Orden lógico\\
\midrule
inicial      & & & & & & \\
N1 (no come) & & & & & & \\
N2 (come)    & & & & & & \\
N3 (no come) & & & & & & \\
N4 (no come) & & & & & & \\
N5 (come)    & & & & & & \\
N6 (no come) & & & & & & \\
\bottomrule
\end{tabular}
\renewcommand{\arraystretch}{1}

% Observaciones por movimiento (resize en N1, envoltura en N3--N4, resize con envoltura en N6).


%=====================================================================
\section*{P4 --- Análisis amortizado}

\sub{a}{Aggregate Method}
% Longitud maxima transitoria n <= g+3. Capacidades en cada resize. Sumar copias en funcion de g.


\sub{b}{Suma geométrica}
% Por que cada termino es la capacidad vigente exactamente cuando ocurre el resize.
% No basta escribir 1+2+4+...=O(n).


\sub{c}{Costo total y amortizado de \texttt{addHead}}
% T(m) y T(m)/m.


\sub{d}{Papel de $m$ y $g$}
% Que parte del costo depende de m y cual de g.


\sub{e}{Método del potencial}
\[
\Phi(\texttt{size},C)=
\]
\begin{itemize}
\item \texttt{addHead} sin resize:
\item \texttt{resize}:
\item \texttt{removeTail}:
\end{itemize}
% Suma de costos amortizados vs. reales (Phi_0, Phi_fin).
% Suficiencia del potencial pese a que removeTail elimina elementos.


%=====================================================================
\section*{P5 --- The Adversarial Snake}

\sub{a}{El adversario no cambia el número de copias}
% Prueba: la secuencia de capacidades depende solo de la politica.
% Que capacidad adicional necesitaria el adversario?


\sub{b--e}{Comparación de políticas}
% Definir: n, inserciones, como se mide el espacio sin usar y el costo amortizado.
\par\smallskip
\renewcommand{\arraystretch}{2.4}
{\setlength{\tabcolsep}{3pt}\begin{tabular}{@{}L{1.45cm}L{3.2cm}L{4.3cm}L{3.2cm}L{3.7cm}@{}}
\toprule
Política & \# Resizes & Total copias & Costo amortizado por crecimiento & Espacio sin usar tras resize\\
\midrule
$C+1$              & & & & \\
$2C$               & & & & \\
$\lceil1.5C\rceil$ & & & & \\
\bottomrule
\end{tabular}}
\renewcommand{\arraystretch}{1}
% Justificacion de cada fila.


\sub{f}{Trade-off}
% Tiempo, frecuencia de resize, copias, memoria reservada.


\medskip\par\noindent\textbf{Familia $C_{\text{new}}=\lceil\alpha C\rceil$.}
\begin{itemize}
\item Valores de $\alpha$:
\item Progreso ($C_{\text{new}}>C$):
\item Cota superior de copias:
\item Cota inferior de copias:
\item Costo amortizado:
\item Necesidad de $\alpha>1$ constante:
\end{itemize}

\end{document}